\section{Introduction}
An instruction to paint ``in the style of Monet'' can introduce a shared
painting appearance or change features that distinguish painters. Moving closer
to Monet's collection does not establish agreement with the Monet--Sisley
difference: a generic painting instruction could improve proximity to both.

We ask whether painter names induce differences aligned with those among
\emph{documented digital reproductions}, beyond a response shared across names.
We center each painter's mean color, spatial and texture measurements at the
four-painter average and evaluate agreement with these labeled contrasts.
Historical subject mixtures and reproduction processes contribute to this finite
target. Fixed generated scene text controls the prompt intervention without
matching historical subjects.

A generator can separate named conditions whose labels map to the wrong
reference painters, or match the reference pattern while adding unrelated
differences. A restrained response can have lower squared error without better
artist distinction. Proximity and recognition alone do not separate these cases.

Six configurations receive the same scenes and clauses. Cross-repeat products
correct random variation under explicit independence assumptions; post-result
checks examine painter pairs, contrast magnitude, scenes and reference quality.

Our contribution is a same-image evaluation study of two distinct common
components. Shared movement across names supplies most own-painter prototype
gain in CLIP and a released CSD encoder. A separate generated-to-reference mean
offset affects prompted-name decisions: an untuned correction fitted on other
scenes changes held-scene accuracy, including adverse CLIP outcomes.
Meanwhile, the apparent failure to distinguish a related painter pair depends
on the representation. These established geometric operations diagnose what
proximity and recognition measure; they introduce no new algorithm or
perceptual criterion.


\section{Related Work}
\citet{somepalli2024style} introduce CSD and General Style Similarity: a generated
embedding's dot product with an averaged artist prototype, including
content-constrained prompts. \citet{su2025prompted} evaluate prompted-artist
recognition using content-controlled name substitutions and matched artist-free
images; generated-image retrieval usually outperforms real-image retrieval in
their benchmark. Our stronger generated-prototype comparator therefore
contextualizes an established domain distinction. We audit common and labeled
contributions to proximity gain, not a new recognition task.

\citet{frochte2026csd} diagnose raw CSD cosine and shared-tradition confusions,
then evaluate local-density readouts and input-resolution changes. Their
absolute-score and pair-verification study excludes closed-set discrimination.
Our controlled name interventions quantify common/labeled gain and connect a
separate mean offset to prompted-name decisions. The general warning that raw
style similarity needs validation is established, not our novelty claim.

\citet{deliege2025} assess stylistic imitation and stereotyping with experts.
AI-Pastiche studies authenticity and prompt adherence \citep{asperti2025pastiche};
AI-WikiArt uses painting-derived captions and vision-language models for
attribution \citep{fu2025aiwikiart}. Our measured painter contrasts and same-image
CLIP/CSD comparisons target neither convincing imitation nor physical authorship.

Quantitative art research examines color, composition and image statistics
\citep{kim2014,sigaki2018,lee2020}. Unlike the contextual image-to-image
experiments of \citet{kim2026}, our text-only intervention supplies no original
composition.

Generative evaluation separates fidelity and diversity \citep{naeem2020},
while conditional metrics distinguish within-condition variation from relationships
among means \citep{benny2021conditional}. We use independent-partition inner products
as the primary noise-corrected squared-difference estimator \citep{diedrichsen2017representational}
and energy statistics as a complementary distributional measure \citep{szekely2013}.
Processing sensitivity \citep{parmar2022} and limits of learned-embedding
separation \citep{asperti2026} motivate the source-quality checks and the
restriction to a specified digital reference target.

Target-domain centroid classifiers also address classifier bias under domain
shift \citep{liang2021atdoc}. Our retrospective decision diagnostic uses a fixed
mean translation and supervised centroids with frozen features; it is not an
adaptation-method contribution.

\section{Design and Measurements}\label{sec:design}
\subsection{Painters and Reference Images}
The primary panel covers four related painters through 649 distinct reproductions: 297 Monet,
106 Sisley, 141 Pissarro and 105 C\'ezanne. Artist means receive equal weight,
so the larger Monet collection does not dominate the between-artist target.
A separate 221-work development panel sets common feature units; no primary
reference or generated image is used to fit that transformation. Both panels
were assembled before the present experiment and had been used in earlier
analyses. They are explicit, finite comparison panels rather than newly withheld
samples of each painter's complete oeuvre.

Reference files carry Wikimedia Commons source metadata and Wikidata work
identities \citep{commonsglam,vrandecic2014wikidata}. These support traceable
selection and rights records, but do not standardize photography, color
calibration or conservation state. These are measurements of digital
reproductions, whose capture conditions can affect the comparison.

\subsection{A Common-scene, Six-model Experiment}
We compare six requested service configurations, labeled GPT Image 1, GPT Image 2, GPT Image 2.5 Flare (Flare),
GPT Image 2.5 Sunburst (Sunburst), Nano Banana 2 and FLUX.2 Max (FLUX).
Fourteen fixed outdoor scene descriptions (briefs) span water,
built, land and mixed compositions. Each scene is crossed with six clauses:
no painting instruction; ``Render as an oil painting''; and ``Render as an oil
painting in the style of [painter]'' for each of the four painters. The scene
text is otherwise identical. Each model--scene--clause cell receives two
separately requested outputs, giving 1,008 images collected on 10 September 2026 (UTC).

All models receive text only and return 1,024-by-1,024-pixel images.
The four OpenAI requests specify medium quality; Nano Banana 2 specifies
1K resolution, while FLUX requests PNG output without an explicit resolution
field. These comparisons concern the recorded
configurations, not maximum attainable quality or equal expenditure per image.
Response records do not echo model identifiers, so these labels identify requested
configurations rather than independently verified checkpoints
(Appendix~\ref{app:provenance}).
The exact prompts, parameters and randomized request order were fixed before
collection. The 14 scenes are a fixed design,
not a random sample of all possible content.
Earlier, separate prompt interventions are summarized in Appendix~\ref{app:cohorts}.

Appendix Figure~\ref{fig:original-generated} shows all six conditions for one brief.
GPT Image 2's C\'ezanne trees have darker blue-green outlines than its Monet
and Sisley trees: a visible name response in this scene, whose agreement with
historical painter differences remains to be measured.



\subsection{What the Features Measure}\label{sec:features}
The 31 features in Table~\ref{tab:features} capture properties that need not
change together. Two images can have similar chroma (color strength) but different
hue diversity, or similar total contrast but different edge directions and
fine-scale texture. Such distinctions make results by feature family interpretable.
There are 11 color coordinates (lightness, chroma, hue and color increments),
eight spatial coordinates (spectra, gradients, orientations and quadrants),
and 12 texture coordinates (wavelet energy, local binary patterns and local contrast).
Correlated coordinates and sensitivity to image processing nevertheless prevent
treating Euclidean distance as a validated measure of artistic style.



Images undergo orientation correction and conversion to sRGB using valid
embedded profiles; compatible unprofiled files are treated as sRGB. The primary
pipeline preserves aspect ratio at a 512-pixel short side, using Lanczos resizing
without upsampling. No painting-region mask is applied: source margins or
calibration strips can enter the measurements. Each coordinate is centered and scaled by its
median and interquartile range in the development panel, with equal total
weight per painter. Thus a unit is a development-panel IQR in that coordinate,
not a perceptual just-noticeable difference. The Euclidean metric weights
standardized coordinates equally, without further equalizing feature families.



\section{Comparing Painter Differences}\label{sec:method}
\subsection{What Counts as Agreement?}\label{sec:contrast-intuition}
The comparison has three steps. Measure the same 31 properties in every image.
For each historical painter, subtract the average of all four painters; this
records how that painter differs from the group. Repeat the subtraction among
the four named outputs for each generated scene, then compare the corresponding
painter differences. We call these differences \emph{contrasts}.

Consider median lightness: Monet averages $.325$ development-IQR units below
Sisley in the references. Across all 14 scenes and both repeats, GPT Image 2
has a smaller gap, $.254$, in the same direction. Subtracting each set's
four-painter mean preserves these gaps; a shared brightness shift cannot create
this distinction. The full comparison uses all 31 coordinates and four painters.

Two scores answer different questions. The \emph{aligned response}, $\beta$,
measures how far generated contrasts extend along the reference contrasts;
$\beta=1$ matches that component's strength. The \emph{error}, $D$, counts total
disagreement, including differences in other directions. Consider three idealized
generators: one makes no painter distinctions, one exactly matches the reference
contrasts, and one doubles them. Their expected errors are respectively 1, 0
and 1. A stronger response can therefore have more error, and a near-unit
aligned response can coexist with substantial mismatch elsewhere.

\subsection{Computing the Two Scores}
Let $\mu_a\in\R^{31}$ be painter $a$'s standardized reference mean and
$r_a=\mu_a-\frac14\sum_b\mu_b$ its contrast. The total squared size of these
contrasts, $H=\sum_a\|r_a\|^2$, normalizes both scores. Here $\|u\|^2$ sums squared
coordinates, $\langle u,v\rangle$ sums matching-coordinate products, and
hats mark estimates. For generated measurements $z_{msak}$, indices denote
model $m$, scene $s$, painter $a$ and repeat $k\in\{1,2\}$:
\begin{equation}
 d_{msak}=z_{msak}-\frac14\sum_b z_{msbk}.
 \label{eq:center}
\end{equation}
Centering removes shared additive shifts, but shared rescaling or mixing
of coordinates can still change the contrasts
(Appendix~\ref{app:score-details}). With $S=14$ scenes, the aligned response is
\begin{equation}
\begin{gathered}
\widehat\beta_{msk}=\frac{\sum_a\langle d_{msak},r_a\rangle}{H},\\
\widehat\beta_m=\frac1S\sum_s\frac{\widehat\beta_{ms1}+\widehat\beta_{ms2}}2.
\end{gathered}
\label{eq:beta}
\end{equation}
A positive value indicates an overall response along the reference differences;
it does not require agreement for each artist pair. Values above one indicate
exaggeration along that direction. To measure error, multiply the discrepancies
from the two repeats rather than squaring either repeat's discrepancy:
\begin{equation}
\begin{gathered}
\widehat D_{ms}=\frac1H\sum_a
\langle d_{msa1}-r_a,\ d_{msa2}-r_a\rangle,\\
\widehat D_m=\frac1S\sum_s\widehat D_{ms}.
\end{gathered}
\label{eq:distortion}
\end{equation}
Under stable conditional means and independent, mean-zero repeat errors,
noise that does not recur across requests contributes zero in expectation.
The expected score is then the normalized squared error of the scene-conditional
mean contrasts. Individual estimates can be negative; they are retained, not
interpreted as better-than-perfect recovery. The $D=1$ no-contrast benchmark
is a theoretical predictor, not the observed generic-painting arm.

The target is the same pooled historical contrast in every scene. A model
whose artist differences vary across scenes can therefore incur error even
if its across-scene average matches the references. Such variation need not
be an artistic mistake. The primary score audits whether the pooled reference
contrast recurs across the tested scenes; it does not identify the historically
appropriate contrast for each generated composition. We separately report mismatch after averaging scenes
and the added scene-variation component; their sum is $D$. Exact identities
and derivations are in Appendix~\ref{app:score-details}.

\subsection{Model Comparisons and Uncertainty}
We compare models scene by scene, then average their paired error differences.
This controls for the common brief; features, painters and individual images
are not treated as independent replicates.

Before inspecting new feature outcomes, we declared six aligned responses and
15 model-pair error differences. Paired-scene Student intervals use Bonferroni
adjustment for these 21 comparisons together, giving model-dependent,
approximate 95\% simultaneous intervals. Intervals excluding zero resolve
their comparisons under that model.
Absolute-$D$ intervals are nominal 95\% summaries: they are unadjusted and descriptive.

We condition on the reference panel, scaling and 14 briefs. Repeated requests
with stable conditional means and independent scene/repeat error blocks define
the sampling model; targets average the conditional contrasts or errors over
these fixed scenes. The scene standard errors also include fixed scene
heterogeneity (Appendix~\ref{app:diagnostics}), so they do not isolate
fresh-request variance. The Student approximation neither guarantees coverage
for the actual service nor supports inference to unseen scenes. Two repeats
cannot verify these assumptions.
A retrospective covariance sensitivity gives a Nano Banana 2--FLUX point-order
crossing at an assumed common trace correlation of $.251$
(Appendix~\ref{app:repeat-covariance}). Actual dependence is unidentified;
these scenario curves do not repair interval coverage.

Declared checks vary scenes, reference samples, feature families and measurement
windows. A separately declared sensitivity equalizes four title-derived subject
classes within each painter. These choices probe dependence on the reference
panel and metric; they do not match historical compositions. Their full
specifications remain in Appendix~\ref{app:estimation}.

\subsection{What the Additional Comparisons Test}\label{sec:diagnostics}
The artist-free and generic-painting arms characterize the common response removed
by centering.
Retrospective diagnostics address three questions: which painter pairs contribute
to the aggregate response, how contrast magnitude changes error, and how reference
choices change the comparison. For the magnitude check, we choose one
nonnegative multiplier per model, shared across all painter contrasts and
coordinates, to minimize error on 13 scenes.
We evaluate it on the omitted scene, rotating through all 14 scenes.
This \emph{calibration} multiplies only centered generated contrasts; reference
contrasts, coordinate units and images stay fixed. These diagnostics
interpret the declared results without extending their significance-test family.
Formulas appear in Appendix~\ref{app:score-details}; controls, fits and the source
audit are in Appendix~\ref{app:diagnostics}.

\section{Results}\label{sec:results}
\subsection{A Response in the Right Direction Can Still Miss the Target}
All 1,008 outputs yield valid measurements. With the full-frame
reference target, every model's aligned-response interval lies above zero
(Table~\ref{tab:specificity-models}). A purely common additive response is
therefore insufficient to explain the named conditions under the stated
inference model. Yet positive alignment does not ensure low error. GPT Image 2
nearly matches the reference-aligned strength, $\beta=.999\;[.893,1.104]$,
while its total error is $D=1.226$: it adds substantial differences beyond
that aligned component.

\begin{table*}[tbp]
\caption{Agreement with the full-frame reference target. $\beta=1$ matches its aligned amplitude; larger is not necessarily better. Expected $D$ is 0 for exact agreement and 1 for no painter distinctions. Positive $\Delta D$ favors FLUX.2 Max. Table brackets give approximate 95\% simultaneous intervals from the prespecified family of six aligned responses and 15 scene-paired model contrasts. FLUX's lowest $D$ estimate has an unadjusted 95\% interval $[0.586,1.016]$, spanning the no-distinction value 1; this interval is descriptive.}
\label{tab:specificity-models}

\centering\small\setlength{\tabcolsep}{4pt}
\begin{tabularx}{\linewidth}{@{}>{\raggedright\arraybackslash}p{.24\linewidth}>{\raggedright\arraybackslash}Xr>{\raggedright\arraybackslash}X@{}}
\toprule
Model & Aligned response $\beta$ & \shortstack{Uncalibrated\\error $D$} & $\Delta D$ versus FLUX\\
\midrule
GPT Image 1 & $0.938\;[0.709,1.168]$ & $1.572$ & $0.771\;[-0.037,1.580]$\\[3pt]
GPT Image 2 & $0.999\;[0.893,1.104]$ & $1.226$ & $0.425\;[-0.161,1.011]$\\[3pt]
GPT Image 2.5 Flare & $0.772\;[0.680,0.864]$ & $1.738$ & $0.937\;[0.256,1.618]$\\[3pt]
GPT Image 2.5 Sunburst & $0.824\;[0.680,0.968]$ & $1.641$ & $0.840\;[0.058,1.623]$\\[3pt]
Nano Banana 2 & $0.442\;[0.256,0.629]$ & $1.109$ & $0.309\;[-0.847,1.464]$\\[3pt]
FLUX.2 Max & $0.470\;[0.239,0.701]$ & $0.801$ & $0$ (baseline)\\[3pt]
\bottomrule\end{tabularx}

\end{table*}


FLUX.2 Max has the lowest error estimate, $.801$, despite a smaller aligned
response, $.470$. Adjusted comparisons resolve higher error for Flare and Sunburst
than for FLUX; the other 13 intervals include zero (Appendix~\ref{app:results},
Table~\ref{tab:all-specificity-pairs}).
These use full-frame references;
the Sunburst result changes after source correction
(Section~\ref{sec:source-correction}). No model is established to beat the
no-contrast benchmark of one. Thus the study detects a
reference-aligned response more clearly than it establishes low total mismatch.

Appendix~\ref{app:geometry} visualizes the same labeled contrasts in two
reference-fitted axes; Figure~\ref{fig:artist-pairs} uses all 31 features.



\subsection{Most Scene-averaged Change Is Common to the Painter Names}\label{sec:shared-change}
After averaging scenes, named-minus-artist-free shifts comprise a four-name mean
and painter-specific departures. The mean includes the oil-painting clause and
any common naming response: 82.5--95.7\% of repeat-corrected squared change
(Table~\ref{tab:shared-change}). Retrospectively changing the baseline to generic
oil painting isolates the additional naming change: 66.7--88.4\% is shared.
The painter-specific component is unchanged. Deleting each scene and recomputing
the ratios gives 65.9--90.2\% across models; these are sensitivity ranges,
not confidence intervals (Appendix~\ref{app:direct-naming}).

\begin{table*}[htbp]
\caption{Scene-averaged, repeat-corrected change. Shared is the four-name mean shift; between-name is departure from it. The artist-free baseline includes the painting clause; the retrospective generic-baseline row isolates additional naming. The last two rows compare the shared named-minus-free shift with generic-minus-free: corrected cosine measures directional agreement, and the ratio divides their corrected squared magnitudes, shared by generic.}
\label{tab:shared-change}
\centering\small\setlength{\tabcolsep}{3pt}
\begin{tabularx}{\linewidth}{@{}l*{6}{>{\centering\arraybackslash}X}@{}}
\toprule
& \shortstack{GPT Image\\1} & \shortstack{GPT Image\\2} & Flare & Sunburst & \shortstack{Nano\\Banana 2} & \shortstack{FLUX.2\\Max}\\\midrule
Free baseline: shared (\%) & 82.5 & 85.8 & 84.8 & 83.9 & 88.1 & 95.7\\
Free baseline: between-name (\%) & 17.5 & 14.2 & 15.2 & 16.1 & 11.9 & 4.3\\
Generic baseline: shared (\%) & 71.3 & 70.9 & 71.1 & 66.7 & 69.1 & 88.4\\
\midrule
Corrected cosine & 0.689 & 0.774 & 0.766 & 0.837 & 0.930 & 0.916\\
Squared-size ratio & 1.986 & 1.960 & 2.816 & 3.356 & 3.692 & 4.138\\
\bottomrule\end{tabularx}

\end{table*}


The shared named-minus-free shift points broadly along the generic oil-painting shift
(noise-corrected cosine $.689$--$.930$; 1 means parallel), but its repeat-corrected
squared magnitude is 1.96--4.14 times as large. The generic and additional naming
shifts interact: their signed cross term is negative for GPT Image 1 and positive
for the other five configurations (Appendix~\ref{app:direct-naming}). Their
squared magnitudes therefore do not add as separate contributions.
Computing the free-baseline decomposition within scenes instead gives
88.6--97.2\% shared change (Appendix~\ref{app:estimation}).


\subsection{Uneven Painter Distinctions in the 31 Features}
Retrospective pair diagnostics extend the lightness example in
Section~\ref{sec:contrast-intuition} to all 31 features (Figure~\ref{fig:artist-pairs}).
Monet--Sisley aligned responses range from $-.249$ to $.129$ in five models,
versus $.402$ for GPT Image 2. Yet GPT Image 1 has lower estimated error for this pair
($.56$ versus $1.16$): stronger alignment does not ensure lower mismatch.
Omitting C\'ezanne lowers FLUX's overall response from $.470$ to $.096$ and
GPT Image 2's from $.999$ to $.651$.

Swapping Monet and Sisley lowers error for Flare, Sunburst and Nano Banana 2.
Reversing a negatively aligned pair difference lowers error without changing its size.

\begin{figure*}[!htbp]
\centering\includegraphics[width=\linewidth]{specificity_artist_pairs.pdf}
\caption{Descriptive painter-pair estimates (no pairwise tests). M, S, P and C denote Monet, Sisley,
Pissarro and C\'ezanne. Pair aligned response 0 means no aligned component,
1 matches reference strength, values above 1 exaggerate it, and negative values
indicate the opposite direction.
Expected error is 0 for exact agreement and 1 for no distinction. Each pair uses
all 31 features and is normalized by its squared reference distance.
Negative error can arise from repeat correction,
not better-than-exact agreement.}
\label{fig:artist-pairs}
\end{figure*}


\subsection{Scene Averaging and Contrast Magnitude Change the Ordering}
The error target matters even within the same coordinates.
Averaging scenes gives Nano Banana 2 lower error than FLUX ($.723$ versus $.761$),
whereas scene-wise scoring reverses that order because Nano Banana 2 varies more
across scenes. Rescaling centered contrasts by a multiplier fitted on the other
13 scenes instead gives GPT Image 2 the lowest held-scene error ($.554$ versus
FLUX's $.714$). These operations change the evaluated feature contrasts, not
images; lower error is not evidence of improved fidelity. The full decomposition,
calibration and influence checks are in Appendix~\ref{app:extended-magnitude}.

\subsection{Which Comparisons Survive Reference Changes?}\label{sec:source-correction}
Within the 31-feature metric, FLUX retains the lowest uncalibrated scene-error point
estimate across scene deletions, measurement windows, reference targets and
feature-weighting sensitivities (Appendices~\ref{app:declared-sensitivities}
and~\ref{app:diagnostics}).

Source handling has a more consequential effect on statistical comparisons.
AI assistants inspected all reference and development reproductions, selecting
crops that exclude peripheral artifacts in 90 reference and 41 development images.
The audit used no human verification. After cropping and refitting development scaling,
FLUX's error estimate changes from $.801$ to $.872$. Corrected aligned-response
intervals remain above zero for all six models, and GPT Image 2 retains the lowest
calibrated estimate.

Appendix Table~\ref{tab:source-comparison} makes the changed comparisons explicit. With
corrected regions and refitted scaling, the Sunburst interval crosses zero,
Flare still has higher error than FLUX and the GPT Image 1 interval moves above zero. The latter
is a retrospective sensitivity, not a new confirmatory finding. Source correction
preserves FLUX's lowest point estimate while changing the evidence for particular
model differences.
Weak Monet--Sisley alignment also persists after correction, but individual
signs are unstable: Flare changes from $-.011$ to $.019$, GPT Image 1 from
$.074$ to $-.053$, and FLUX from $.129$ to $.019$. The label-swap observation
above therefore concerns the original reference target.






A separate check uses class-specific reference means for 11 water, built and
land briefs, excluding three mixed briefs. Replacing title-derived reference
labels with visual labels retains FLUX's lowest error estimate.
Appendix~\ref{app:diagnostics} gives the controls and audit.\par



\subsection{What Changes in Learned Representations?}\label{sec:learned-results}
We remeasure all original images using CLIP ViT-L/14 \citep{radford2021}
and the public CSD ViT-L style checkpoint \citep{somepalli2024style}.
This retrospective protocol was fixed before learned outcomes. Each encoder
uses its native 224-pixel center crop and 768-dimensional unit vectors, without
IQR scaling. This changes representation and native preprocessing together;
the hand-feature square-window sensitivity in Appendix~\ref{app:direct-naming}
does not isolate encoder effects. Reference prototypes are means of the same 649 works;
Appendix~\ref{app:learned-audit} specifies checkpoints, checks and sensitivities.
CSD's release warns of a discrepancy from published benchmark results, so
we evaluate this released checkpoint without claiming benchmark reproduction.

\emph{Common movement also dominates prototype gain.}
Let $g_a$ and $\mu_a$ be mean named and reference embeddings, $g_0$ the generic
mean, and bars denote equal-painter averages. Using unnormalized prototypes,
mean own-painter cosine gain decomposes exactly as
\begin{equation}
\frac14\sum_a g_a^\top\mu_a-g_0^\top\bar\mu
=(\bar g-g_0)^\top\bar\mu+\frac{H\widehat\beta}{4}.
\label{eq:prototype-main}
\end{equation}
The first term is common movement; the second is labeled agreement.
Across configurations, the common term supplies 73.2--83.8\% of CLIP gain and
54.2--79.7\% of CSD gain (Table~\ref{tab:learned-main}). These are contributions
to cosine gain, distinct from the squared-change shares in
Table~\ref{tab:shared-change}. The latter are also majority-common:
76.1--82.3\% in CLIP and 68.8--77.4\% in CSD for naming beyond generic painting.

\input{icml_learned_table.tex}

\emph{Monet--Sisley alignment depends on representation.}
Monet--Sisley alignment is positive for all six configurations in both learned
representations: $.299$--$.532$ in CLIP and $.323$--$.587$ in CSD. Thus the
weak or reversed response measured by the 31 features cannot establish a
representation-independent failure to distinguish these painters. CLIP gives
GPT Image 1 the lowest scene-error estimate ($.847$); CSD gives FLUX the
lowest ($.714$). Distances across representations do not share a perceptual scale.

\emph{Proximity and recognition answer different questions.}
Classifying each named output by its closest unit-normalized reference prototype
gives the accuracies in Table~\ref{tab:learned-main}. Nano Banana 2 has the largest
CLIP prototype gain ($.1121$) but 51.8\% accuracy, versus GPT Image 2's smaller
gain ($.0996$) and 68.8\% accuracy. In CSD, GPT Image 2 has the smallest gain
($.1659$) and highest accuracy (75.9\%). These descriptive reversals show why
unpartitioned proximity gain is insufficient for evaluating labeled response.

Classifying the 221 development works with these primary prototypes yields
79.8\% macro accuracy in CLIP and 79.6\% in CSD. Appendix~\ref{app:learned-audit}
reports their painter confusions, reference energies and all absolute
common/labeled gain components. These controls describe panel separability,
not generated-image fidelity.

The appendix reports every configuration under both representations, original
and audited regions, and primary and development targets. These related
encoders share a CLIP backbone and unknown training-image overlap; the same
images supply all comparisons. They add a representation sensitivity, not a
fresh service replication or perceptual validation. No new significance claims
are attached to this audit.

\subsection{Common Offsets Can Change Prompted-name Decisions}\label{sec:transfer}
A decision task tests a consequence beyond score disagreement: recovering the
known painter clause of a generated image. This retrospective extension was
chosen after inspecting baseline confusions; its rules were fixed before
computing corrected outcomes. It uses all configurations and both encoders.
For each held-out scene, pool the 104 named outputs from the other 13 scenes,
excluding both held-scene repeats. Their mean $\bar g_{-s}$ ignores individual
painter assignments within the known balanced pool. With raw real-art centroids
$\mu_a$ and normalized directions $u_a=\mu_a/\|\mu_a\|$, set
\begin{equation}
 t_{-s}=\tfrac14\sum_a\mu_a-\bar g_{-s},\qquad
 \hat a=\arg\max_a (x+t_{-s})^\top u_a.
 \label{eq:transfer-main}
\end{equation}
The translation scale is fixed at one. It aligns mixture means and changes
class intercepts by $t_{-s}^\top u_a$; it leaves centered painter differences
unchanged within each fold. This cross-domain offset is not the
named-minus-generic effect in Equation~\ref{eq:prototype-main}: free and generic
outputs never enter its fit.

\input{icml_transfer_table.tex}

On the original-view primary target, translation changes accuracy by $-5.4$
to $+12.5$ percentage points in CLIP,
with two declines, one tie and three gains; all six CSD changes are positive,
from $+0.9$ to $+26.8$ points. Equal-configuration means are $+2.7$ and $+10.0$
points, respectively. These summaries retain newly introduced errors:
CSD FLUX corrects 39 decisions and harms nine of 112, while CLIP GPT Image 1
corrects one and harms seven. Thus a common offset can affect recognition
without changing painter contrasts, but mean matching is not uniformly useful.

A supervised context rule fits one centroid per prompted painter on those same
13 training scenes. Its accuracy is higher than both reference rules in all 12
combinations (Table~\ref{tab:transfer-main}); it uses artist assignments that
translation does not. This demonstrates recoverable prompt-label structure
within generated images, not agreement with historical painting style.
Appendix~\ref{app:transfer} reports all confusions, painter recalls, scene changes
and reference sensitivities. Folds overlap and reuse previously inspected
images; these are descriptive out-of-fold results, not independent validation.

\section{Discussion}\label{sec:discussion}
Detecting a name response and matching painter differences are distinct achievements.
FLUX's lowest uncalibrated error in the 31-feature metric coexists with weak
Monet--Sisley alignment, while learned representations yield positive pair alignment.
Model comparisons should therefore report overall and pairwise scores together
for the stated reference target, distinguishing estimates from resolved comparisons.
Reports should also state whether they score each scene, the average pattern,
or rescaled contrasts. For prompted-name recognition, the two held-scene
controls distinguish sensitivity to a pooled domain offset from supervised
within-generator separability. Neither lower contrast error nor better name
recovery alone establishes closer painter resemblance.

The intervention identifies differences between prompt conditions, not their
internal causes. Training composition, guidance and shared transformations remain
possible explanations.

\subsection{Scope and Remaining Limitations}
The target is the pattern of differences in a finite digital reference
collection. It mixes subject choices and reproduction processes with properties
of the paintings. Coarse content classes do not match composition, viewpoint,
career period or actual generated scene adherence. The source audit addresses
visible peripheral regions and title/visual-class disagreements, but AI
judgments, uncertain boundaries and remaining capture differences limit that
correction. One reproduction per work cannot separate effects of illumination,
sharpening, conservation or digitization source.

The 31 coordinates describe digital color, spatial organization and texture;
neither they nor the learned encoders establish physical brushwork or perceived
painter resemblance. CSD's style-tag list includes all four painters, and exact
training-image overlap is unknown for both encoders. Numerical replay and
representation sensitivities do not validate perceptual judgments.
Inference also conditions on four related
painters, 14 authored briefs, two repeats and one clause template. Simulations
probe stated noise models, not actual service independence. The image panels
make the comparison inspectable, while incomplete public exact-pixel access
still limits independent measurement reproduction.

\section{Conclusion}
Most squared feature change from artist-free to named prompts is shared across
names after scene averaging, including the oil-painting contribution.
This remains true for additional naming relative to generic painting in a
retrospective comparison.
Most learned prototype gain also comes from common movement. Weak or reversed
31-feature Monet--Sisley alignment in five configurations contrasts with positive
point estimates in all six configurations in CLIP and CSD.
Held-scene mean translation changes name-recovery decisions while preserving
centered painter contrasts; its mixed CLIP outcomes prohibit treating it as a
universal correction. Evaluation should separate shared proximity gain,
reference-based recognition and generated-domain separability. None alone
supports a general ranking of painter fidelity.
